\noindent
Zur Berechnung der im Lastwiderstand umgesetzten Leistung wird zunächst der Strom
durch den Lastwiderstand bestimmt. Aufgrund der Parallelschaltung von Innen- und
Lastwiderstand ergibt sich dieser aus dem Stromteiler zu:
\begin{equation}
I = I_0 \frac{G_L}{G_L + G_I}
\end{equation}

\noindent
Die am Lastwiderstand anliegende Spannung ergibt sich aus dem Zusammenhang zwi-
schen Strom und Leitwert:
\begin{equation}
U = \frac{I}{G_L}
\end{equation}

\noindent 
Die elektrische Leistung ergibt sich allgemein aus Spannung und Strom:
\begin{equation}
P = U \cdot I
\end{equation}

\noindent
Durch Einsetzen der Gleichungen für Strom und Spannung in die Leistungsformel er-
gibt sich die im Lastwiderstand umgesetzte Leistung zu:
\begin{equation}
P(G_L) = I_0^2 \frac{G_L}{(G_L + G_I)^2}
\end{equation}